\clearpage
\section{Заключение}

По результатам проведенного исследования
можно сформулировать следующие выводы:

Во-первый, на текущий момент на рынке отсутствует
оптимальное решение для автоматизированного создания
ГЗ в Obsidian.
Представленные инструменты имеют компромиссы.
Корпоративные графовые базы данных (Neo4j и MemGraph)
требуют дополнительную инфраструктуру,
чрезмерно зависят от LLM и не поддерживают инкрементальные обновления.
Облачные сервисы для ведения заметок имеют риски утечки данных
и привязывают пользователей к экосистеме.
Эти недостатки подтверждают необходимость создания
легковесного и независимого от поставщиков алгоритма,
сохраняющего приватность пользовательских данных.

Во-вторых, анализ готовых решений позволил
спроектировать и реализовать
эффективный, отказоустойчивый конвейер
для автоматизированного построения ГЗ в Obsidian.
Разработанное решение совмещает гибридный поиск
(BM25 и семантический по плотным эмбеддингам BGE-M3 на базе LanceDB)
с фильтрацией через Cross Encoder (Jina)
и типизацию связей с помощью LLM (Gemma-4-31B).
Алгоритм также имеет модель для сохранения состояния конвейера
(\texttt{MetadataManager}),
позволяющий возобновлять работу алгоритма
с того же этапа,
и состояния хранилища (\texttt{VaultManager}),
проверяющий заметки на наличие изменений
посредством сравнения сохраненного списка файлов
с предыдущего запуска и SHA-256 хеша их содержимого,
и обеспечивающее инкрементальные обновления.
Такой подход позволяет
экономить вычислительные ресурсы,
что крайне важно для индивидуального использования.

В-третьих, для тестирования собственного решения,
в частности для сравнения типизированного GraphRAG и нетипизированным,
был создан модифицированный алгоритм GraphRAG,
основанный на решении ObsidianRAG.
Его особенностями стали учет типов связей при сборе контекста
многошаговый обход графа.

В-четвертых, для оценки системы и сравнения с аналогами
была разработана методология тестирования и
собран валидационный набор данных.
Выбор Википедии как основного источника данных для заметок
позволил соблюсти строгие ограничения касательно
содержимого заметок и структуры хранилища.
Собранный датасет состоит из 63 заметок, имеющих 246 связей,
и разделен на 5 сценариев, каждый из которых
проверяет работоспособность алгоритма построения ГЗ
в одном из сложных случаев,
таких как омонимия, пересекающиеся темы, многошаговая хронология и
недостаток контекста.

В-пятых, был проведен ряд сравнений и экспериментов по улучшения
полученного решения.
Полученные результаты подтвердили высокую эффективность и
точность работы предложенных алгоритмов.
Выявлено, что оптимальной стратегией извлечения контекста
является гибридный поиск: она обеспечивает наилучший баланс,
достигая значения F1-меры 0.680.
Для типизации связей лучше всего использовать
few-shot промпт с детально описанными допустимыми типами и задачей.
Такой подход позволил получить значение взвешенной F1-меры
равным 0.669 (и 0.802 без учета типизации).
Оценка собственного алгоритма GraphRAG
доказала целесообразность улучшений, продемонстрировав
значительный рост метрики Context Recall, получив 0.96 в среднем
и 1.0 на подкорпусах \texttt{hub} и \texttt{personal} в случае
нетипизированного обхода и 0.940 в среднем для типизированного,
при полном отсутствии галлюцинаций
(Answer Faithfulness была не ниже 0.88).
Для сравнения, ObsidianRAG имеет среднее значения CR равное 0.786,
а AC = 0.496.
Важным достижением стало минимальное изменение
полученных значений при переходе от размеченного вручную ГЗ
к сгенерированному: падение CR составило 0.013, AC~--- 0.01
в среднем по всем алгоритмам.
Также было проведено сравнение качества извлечения на сгенерированном
с помощью neo4j Graph Builder.
Это решения имеет наивысшие значения
CR и AC равные 0.977 и 0.602 соответственно.

Разработанный алгоритмов является
практически применимым решением,
устраняющим когнитивное «узкое горлышко»,
заключающееся в ручном связывании многочисленных заметок
в системах управления знаниями.
Конвейер позволяет без участия человека превратить
хранилище в граф знаний, который можно использовать
для эффективного извлечения данных с помощью GraphRAG.
Решение демонстрирует качество работы наравне с
продвинутым решением для корпоративного использования (neo4j).

Несмотря на достижения отличных результатов,
в работе есть ряд ограничений.
Прежде всего это ограниченный объем датасета (63 заметки),
однотипность его содержания и узкая онтология типов связей.
Также недостатком можно считать отсутствие дополнительных
шагов, увеличивающих эффективность извлечения, таких как
создание сообществ.

Дальнейшее развитие решения заключается в
проведении оценки на хранилищах большего размера,
состоящих из сотен или тысяч реальных обезличенных
заметок различных форматов и тематик.
Кроме того, перспективно добавить
генерацию сообществ и автоматическую адаптацию
используемых типов связей (объединение, добавление и удаление типов).
Развитие разработанной системы позволит
создать полностью автономный алгоритм
конструирования ГЗ в Obsidian.


% Представленная выпускная квалификационная работа посвящена
% актуальной проблеме преодоления «когнитивного узкого места»
% в системах управления персональными знаниями (PKM). В условиях
% экспоненциального роста объемов информации классические методы
% ведения заметок и ручного связывания данных становятся неэффективными,
% что ведет к деградации персональных баз знаний. В ходе
% проведенного исследования была успешно достигнута главная
% цель работы: разработан, реализован и всесторонне оценен
% автоматизированный алгоритм семантического связывания
% текстовых данных для экосистемы Obsidian, а также
% усовершенствован конвейер генерации, дополненной поиском по графу знаний (GraphRAG).

% На основании проведенного теоретического и практического
% исследования можно сформулировать следующие основные выводы:

% Во-первых, анализ существующих решений показал, что на
% рынке отсутствует оптимальный инструмент для локальных
% баз знаний, который сочетал бы в себе вычислительную
% эффективность, суверенность данных и устойчивость к шуму.
% Корпоративные графовые базы данных (такие как Neo4j) требуют
% дорогостоящей инфраструктуры и не поддерживают инкрементальные
% обновления без перестройки всего графа, а облачные сервисы создают
% неприемлемые риски утечки приватной информации. Это подтвердило
% острую научно-практическую необходимость в создании собственного
% легковесного алгоритма, независимого от закрытых вендорских экосистем.

% Во-вторых, в рамках практической части был спроектирован и
% реализован отказоустойчивый конвейер построения графа знаний.
% Разработанный алгоритм комбинирует гибридный поиск (
% BM25 и плотные BGE-M3 эмбеддинги на базе LanceDB)
% с фильтрацией посредством Cross-Encoder (Jina) и
% типизацией связей с помощью локальных языковых моделей.
% Ключевым архитектурным достижением стала реализация модуля
% управления состоянием (VaultManager и MetadataManager),
% который сканирует хранилище по хешам SHA-256 и обеспечивает
% инкрементальное обновление графа. Это позволяет возобновлять
% работу алгоритма без повторной траты вычислительных ресурсов,
% что критически важно для потребительского оборудования.

% В-третьих, для интеграции построенного графа в процессы поиска
% ответов был разработан усовершенствованный четырехэтапный алгоритм
% GraphRAG. В отличие от референсного решения (ObsidianRAG),
% предложенный пайплайн включает этапы многошагового обхода соседей
% и строгого отбора контекста (Pruning) с переранжированием. Были
% реализованы два режима обхода: базовый (доминирующий на разреженных
% графах) и типизированный (эффективно фильтрующий шум на плотно
% связанных графах).

% В-четвертых, для объективной валидации системы была разработана
% оригинальная методология тестирования и собран репрезентативный
% набор данных на базе статей Википедии. Отказ от сырых пользовательских
% заметок в пользу энциклопедических данных позволил соблюсти
% строгое правило «одна заметка — одна сущность». Датасет из 63
% заметок и 246 связей, разделенный на пять топологических сценариев
% (звездообразный граф, омонимия, пересекающиеся темы, хронология и
% персональные записи), обеспечил комплексную проверку системы на
% способность разрешать смысловые конфликты и находить многошаговые
% логические связи.

% В-пятых, результаты проведенных экспериментов со строго зафиксированными
% параметрами воспроизводимости эмпирически подтвердили высокую
% эффективность предложенных решений. Выявлено, что гибридная
% стратегия поиска кандидатов (Combined) обеспечивает наилучший
% баланс, достигая значения F1-меры 0.680. На этапе классификации
% связей применение few-shot промптинга позволило модели Gemma-4-31B
% превзойти более ресурсоемкие аналоги, показав взвешенную F1-меру
% 0.669 (и 0.802 без учета типизации). Оценка конвейера GraphRAG
% продемонстрировала способность системы извлекать практически
% исчерпывающий контекст (метрика Context Recall достигла 0.96)
% при полном отсутствии галлюцинаций (Answer Faithfulness стабильно
% превышала 0.88). Важнейшим результатом стало то, что переход от
% использования эталонного (размеченного вручную) графа к графу,
% сгенерированному нашим алгоритмом автоматически, привел к падению
% корректности ответа (Answer Correctness) в среднем менее чем на 0.01.
% Это доказывает высочайшую устойчивость разработанной системы к информационному шуму.

% Общий вывод по работе. Разработанный комплекс алгоритмов представляет
% собой законченное, научно обоснованное и практически применимое решение,
% ликвидирующее «когнитивное узкое место» в системах управления персональными
% знаниями. Алгоритм автоматизирует рутинные процессы семантического связывания,
% сохраняя полный суверенитет пользовательских данных и высокую скорость
% работы на локальных машинах. Использование размеченного графа в связке
% с модифицированным многошаговым пайплайном GraphRAG кратно повышает
% полноту и точность извлечения информации по сравнению с традиционными
% методами гибридного поиска.

% Несмотря на высокие достигнутые показатели, исследование имеет ряд
% объективных ограничений, задающих векторы для дальнейшей научной работы.
% В первую очередь, это ограниченный объем (63 заметки) и однотипность
% валидационного датасета, а также узкая онтология типов связей.

% Перспективы дальнейшего изучения темы лежат в плоскости масштабирования
% системы. Планируется проведение нагрузочного тестирования конвейера на
% хранилищах, содержащих десятки тысяч файлов, с привлечением реальных
% обезличенных пользовательских баз. Кроме того, перспективным направлением
% является расширение и динамическая адаптация словаря онтологий
% (типов связей) под специфику конкретного пользователя, а также внедрение
% механизмов кластеризации LLM-предсказаний для устранения путаницы между
% семантически близкими классами отношений. Развитие предложенных методов
% позволит создать полностью автономного «умного ассистента» для любых локальных
% баз знаний, функционирующего без вмешательства пользователя.